\nwfilename{nw/classical.nw}\nwbegindocs{0}% -*- mode: poly-noweb; noweb-code-mode: sml-mode; -*-% ===> this file was generated automatically by noweave --- better not edit it
%%
%% classical.nw
%% 
%% Made by Alex Nelson <pqnelson@gmail.com>
%% Login   <alex@lisp>
%% 
%% Started on  2026-04-21T12:45:57-0700
%% Last update 2026-04-21T12:45:57-0700
%% 

\section{Classical tableaux}

\begin{node}
We want to investigate if certain $FS_{0}$ formulas are theorems, had
we used Classical logic instead of Intuitionistic logic. After all,
every Intuitionistically valid formula is also valid Classically (but
not all Classically valid formulas are Intuitionistically acceptable).
Tableaux is one way to investigate this question.
\end{node}

\nwenddocs{}\nwbegincode{1}\sublabel{NW1mR7zJ-3NMRe1-1}\nwmargintag{{\nwtagstyle{}\subpageref{NW1mR7zJ-3NMRe1-1}}}\moddef{Classical.sig~{\nwtagstyle{}\subpageref{NW1mR7zJ-3NMRe1-1}}}\endmoddef\nwstartdeflinemarkup\nwenddeflinemarkup
signature Classical = sig
  val nnf : Formula.t -> Formula.t;
  val askolemize : Formula.t -> Formula.t;
  val tab : Formula.t -> int -> int;
end;
\nwnotused{Classical.sig}\nwidentuses{\\{{\nwixident{Formula.t}}{Formula.t}}}\nwindexuse{\nwixident{Formula.t}}{Formula.t}{NW1mR7zJ-3NMRe1-1}\nwendcode{}\nwbegindocs{2}\nwdocspar

\nwenddocs{}\nwbegincode{3}\sublabel{NW1mR7zJ-l5jX3-1}\nwmargintag{{\nwtagstyle{}\subpageref{NW1mR7zJ-l5jX3-1}}}\moddef{Classical.sml~{\nwtagstyle{}\subpageref{NW1mR7zJ-l5jX3-1}}}\endmoddef\nwstartdeflinemarkup\nwenddeflinemarkup
structure Classical :> Classical = struct
\LA{}Transform formula to negation normal form~{\nwtagstyle{}\subpageref{NW1mR7zJ-4TkWUc-1}}\RA{}

\LA{}Simplify formulas using classical tautologies~{\nwtagstyle{}\subpageref{NW1mR7zJ-2RorhN-1}}\RA{}

\LA{}Skolemization~{\nwtagstyle{}\subpageref{NW1mR7zJ-4QrsTb-1}}\RA{}

\LA{}Tableaux refuter~{\nwtagstyle{}\subpageref{NW1mR7zJ-yCIr5-1}}\RA{}
end;
\nwnotused{Classical.sml}\nwendcode{}\nwbegindocs{4}\nwdocspar

\begin{node}
Signed formulas $T\phi$ and $F\phi$ encode whether we have to prove
$\phi$ is true or false, respectively. Tableaux constructs a tree from
signed formulas. For example $T(A\land B)$ can extend the current
branch of the Tableaux with two new signed formulas $TA$ and $TB$. But
$F(A\land B)$ creates two new branches in the tableaux, one starts
with $FA$ and the other starts with $FB$.

Continuing in this manner, we continue writing a branch until a
contradiction has been found. More precisely, we have $TA$ and $FA$
occur within the same branch. In this situation, the branch ``closes''
and we write $\times$ to indicate this. What this means is that we
can stop working on this branch, and we can go to the next branch. If
there are no more branches to work on, then the Tableaux refutes the formula.
\end{node}


\begin{definition}
A \define{Literal} is a formula such that it is either an atomic formula
(a predicate, or equality between terms [if working in first-order logic
  with equality])
or the negation of an atomic formula.
\end{definition}

\begin{definition}
A formula $\Phi$ is said to be in \define{Negation Normal Form} if
\begin{enumerate}
\item $\Phi$ is a literal, or
\item $\Phi=A\land B$ where $A$ and $B$ are in negation normal form,
  or
\item $\Phi=A\lor B$ where $A$ and $B$ are in negation normal form, or
\item $\Phi=\forall x\ldotp A$ or $\Phi=\exists x\ldotp A$, where $A$ is again in negation normal form.
\end{enumerate}
In Classical logic, we can always transform \emph{any} formula to be
in negation normal form.
\end{definition}

\begin{node}
The basic gambit for classical tableaux is:
\begin{enumerate}
\item Given a formula $\phi$, transform it to negation normal form $\psi$;
\item Skolemize $\psi$ to obtain $\chi$;
\item Perform the tableaux method on $\chi$.
\end{enumerate}
\end{node}

\begin{node}%163
We transform a formula into negation normal form by replacing
$A\implies B$ by $(\neg A)\lor B$ and so on. We begin by transforming
the basic connectives and quantifiers by applying the transformation
to the immediate subformulas, using the rules:
\begin{enumerate}
\item $NNF(A\land B)=NNF(A)\land NNF(B)$
\item $NNF(A\lor B)=NNF(A)\lor NNF(B)$
\item $NNF(A\implies B)=NNF(\neg A)\lor NNF(B)$
\item $NNF(A\iff B)=(NNF(A)\land NNF(B))\lor(NNF(\neg A)\land NNF(\neg B))$
\item $NNF(\exists x\ldotp A)=\exists x\ldotp NNF(A)$
\item $NNF(\forall x\ldotp A)=\forall x\ldotp NNF(A)$
\end{enumerate}
\end{node}

\nwenddocs{}\nwbegincode{5}\sublabel{NW1mR7zJ-4TkWUc-1}\nwmargintag{{\nwtagstyle{}\subpageref{NW1mR7zJ-4TkWUc-1}}}\moddef{Transform formula to negation normal form~{\nwtagstyle{}\subpageref{NW1mR7zJ-4TkWUc-1}}}\endmoddef\nwstartdeflinemarkup\nwusesondefline{\\{NW1mR7zJ-l5jX3-1}}\nwenddeflinemarkup
(* nnf : Formula.t -> Formula.t *)
fun nnf fm =
  if Formula.is_and fm
  then let val (A,B) = Formula.dest_and fm
       in Formula.mk_and(nnf A, nnf B) end
  else if Formula.is_or fm
  then let val (A,B) = Formula.dest_or fm
       in Formula.mk_or(nnf A, nnf B) end
  else if Formula.is_imp fm
  then let val (A,B) = Formula.dest_imp fm
       in Formula.mk_or(nnf (Formula.mk_not A),
                        nnf B) end
  else if Formula.is_iff fm
  then let val (A,B) = Formula.dest_iff fm
       in Formula.mk_or(Formula.mk_and(nnf A, nnf B),
                        Formula.mk_and(nnf(Formula.mk_not(A)), nnf(Formula.mk_not(B)))) end
  else if Formula.is_forall fm
  then let val(x,A) = Formula.dest_forall fm
       in Formula.mk_forall(x, nnf A) end
  else if Formula.is_exists fm
  then let val(x,A) = Formula.dest_exists fm
       in Formula.mk_exists(x, nnf A) end
  \LA{}Negation normal form for negated formulas~{\nwtagstyle{}\subpageref{NW1mR7zJ-4HwyeF-1}}\RA{}

\nwused{\\{NW1mR7zJ-l5jX3-1}}\nwidentuses{\\{{\nwixident{Formula.dest{\_}and}}{Formula.dest:unand}}\\{{\nwixident{Formula.dest{\_}iff}}{Formula.dest:uniff}}\\{{\nwixident{Formula.dest{\_}imp}}{Formula.dest:unimp}}\\{{\nwixident{Formula.dest{\_}or}}{Formula.dest:unor}}\\{{\nwixident{Formula.is{\_}and}}{Formula.is:unand}}\\{{\nwixident{Formula.is{\_}iff}}{Formula.is:uniff}}\\{{\nwixident{Formula.is{\_}imp}}{Formula.is:unimp}}\\{{\nwixident{Formula.is{\_}or}}{Formula.is:unor}}\\{{\nwixident{Formula.mk{\_}and}}{Formula.mk:unand}}\\{{\nwixident{Formula.mk{\_}exists}}{Formula.mk:unexists}}\\{{\nwixident{Formula.mk{\_}forall}}{Formula.mk:unforall}}\\{{\nwixident{Formula.mk{\_}not}}{Formula.mk:unnot}}\\{{\nwixident{Formula.mk{\_}or}}{Formula.mk:unor}}\\{{\nwixident{Formula.t}}{Formula.t}}}\nwindexuse{\nwixident{Formula.dest{\_}and}}{Formula.dest:unand}{NW1mR7zJ-4TkWUc-1}\nwindexuse{\nwixident{Formula.dest{\_}iff}}{Formula.dest:uniff}{NW1mR7zJ-4TkWUc-1}\nwindexuse{\nwixident{Formula.dest{\_}imp}}{Formula.dest:unimp}{NW1mR7zJ-4TkWUc-1}\nwindexuse{\nwixident{Formula.dest{\_}or}}{Formula.dest:unor}{NW1mR7zJ-4TkWUc-1}\nwindexuse{\nwixident{Formula.is{\_}and}}{Formula.is:unand}{NW1mR7zJ-4TkWUc-1}\nwindexuse{\nwixident{Formula.is{\_}iff}}{Formula.is:uniff}{NW1mR7zJ-4TkWUc-1}\nwindexuse{\nwixident{Formula.is{\_}imp}}{Formula.is:unimp}{NW1mR7zJ-4TkWUc-1}\nwindexuse{\nwixident{Formula.is{\_}or}}{Formula.is:unor}{NW1mR7zJ-4TkWUc-1}\nwindexuse{\nwixident{Formula.mk{\_}and}}{Formula.mk:unand}{NW1mR7zJ-4TkWUc-1}\nwindexuse{\nwixident{Formula.mk{\_}exists}}{Formula.mk:unexists}{NW1mR7zJ-4TkWUc-1}\nwindexuse{\nwixident{Formula.mk{\_}forall}}{Formula.mk:unforall}{NW1mR7zJ-4TkWUc-1}\nwindexuse{\nwixident{Formula.mk{\_}not}}{Formula.mk:unnot}{NW1mR7zJ-4TkWUc-1}\nwindexuse{\nwixident{Formula.mk{\_}or}}{Formula.mk:unor}{NW1mR7zJ-4TkWUc-1}\nwindexuse{\nwixident{Formula.t}}{Formula.t}{NW1mR7zJ-4TkWUc-1}\nwendcode{}\nwbegindocs{6}\nwdocspar

\begin{node}
The rules for transforming a negated formula into negation normal form
are:
\begin{enumerate}
\item $NNF(\neg\neg A)=NNF(A)$
\item $NNF(\neg(A\land B))=NNF(\neg A)\lor NNF(\neg B)$
\item $NNF(\neg(A\lor B))=NNF(\neg A)\land NNF(\neg B)$
\item $NNF(\neg(A\implies B))=NNF(A)\land NNF(\neg B)$
\item $NNF(\neg(A\iff B))=(NNF(A)\land NNF(\neg B))\lor(NNF(\neg A)\land NNF(B))$
\item $NNF(\neg\exists x\ldotp A)=\forall x\ldotp NNF(\neg A)$
\item $NNF(\neg\forall x\ldotp A)=\exists x\ldotp NNF(\neg A)$
\end{enumerate}
\end{node}

\nwenddocs{}\nwbegincode{7}\sublabel{NW1mR7zJ-4HwyeF-1}\nwmargintag{{\nwtagstyle{}\subpageref{NW1mR7zJ-4HwyeF-1}}}\moddef{Negation normal form for negated formulas~{\nwtagstyle{}\subpageref{NW1mR7zJ-4HwyeF-1}}}\endmoddef\nwstartdeflinemarkup\nwusesondefline{\\{NW1mR7zJ-4TkWUc-1}}\nwenddeflinemarkup
else if Formula.is_not fm
then let val fm' = Formula.dest_not fm
     in if Formula.is_not fm'
        then let val A = Formula.dest_not fm'
             in nnf A end
        else if Formula.is_and fm'
        then let val (A,B) = Formula.dest_and fm'
             in Formula.mk_or(nnf(Formula.mk_not(A)),
                              nnf(Formula.mk_not(B))) end
        else if Formula.is_or fm'
        then let val (A,B) = Formula.dest_or fm'
             in Formula.mk_and(nnf(Formula.mk_not(A)),
                              nnf(Formula.mk_not(B))) end
        else if Formula.is_imp fm'
        then let val (A,B) = Formula.dest_imp fm'
             in Formula.mk_and(nnf A,
                               nnf(Formula.mk_not(B))) end
        else if Formula.is_iff fm'
        then let val (A,B) = Formula.dest_iff fm'
             in Formula.mk_or(Formula.mk_and(nnf A,
                                             nnf(Formula.mk_not(B))),
                              Formula.mk_and(nnf(Formula.mk_not(A)),
                                             nnf B))
             end
        else if Formula.is_exists fm'
        then let val (x,A) = Formula.dest_exists fm'
             in Formula.mk_forall(x, nnf(Formula.mk_not(A))) end
        else if Formula.is_forall fm'
        then let val (x,A) = Formula.dest_forall fm'
             in Formula.mk_exists(x, nnf(Formula.mk_not(A))) end
        else fm
     end
else fm;
\nwused{\\{NW1mR7zJ-4TkWUc-1}}\nwidentuses{\\{{\nwixident{Formula.dest{\_}and}}{Formula.dest:unand}}\\{{\nwixident{Formula.dest{\_}iff}}{Formula.dest:uniff}}\\{{\nwixident{Formula.dest{\_}imp}}{Formula.dest:unimp}}\\{{\nwixident{Formula.dest{\_}not}}{Formula.dest:unnot}}\\{{\nwixident{Formula.dest{\_}or}}{Formula.dest:unor}}\\{{\nwixident{Formula.is{\_}and}}{Formula.is:unand}}\\{{\nwixident{Formula.is{\_}iff}}{Formula.is:uniff}}\\{{\nwixident{Formula.is{\_}imp}}{Formula.is:unimp}}\\{{\nwixident{Formula.is{\_}not}}{Formula.is:unnot}}\\{{\nwixident{Formula.is{\_}or}}{Formula.is:unor}}\\{{\nwixident{Formula.mk{\_}and}}{Formula.mk:unand}}\\{{\nwixident{Formula.mk{\_}exists}}{Formula.mk:unexists}}\\{{\nwixident{Formula.mk{\_}forall}}{Formula.mk:unforall}}\\{{\nwixident{Formula.mk{\_}not}}{Formula.mk:unnot}}\\{{\nwixident{Formula.mk{\_}or}}{Formula.mk:unor}}}\nwindexuse{\nwixident{Formula.dest{\_}and}}{Formula.dest:unand}{NW1mR7zJ-4HwyeF-1}\nwindexuse{\nwixident{Formula.dest{\_}iff}}{Formula.dest:uniff}{NW1mR7zJ-4HwyeF-1}\nwindexuse{\nwixident{Formula.dest{\_}imp}}{Formula.dest:unimp}{NW1mR7zJ-4HwyeF-1}\nwindexuse{\nwixident{Formula.dest{\_}not}}{Formula.dest:unnot}{NW1mR7zJ-4HwyeF-1}\nwindexuse{\nwixident{Formula.dest{\_}or}}{Formula.dest:unor}{NW1mR7zJ-4HwyeF-1}\nwindexuse{\nwixident{Formula.is{\_}and}}{Formula.is:unand}{NW1mR7zJ-4HwyeF-1}\nwindexuse{\nwixident{Formula.is{\_}iff}}{Formula.is:uniff}{NW1mR7zJ-4HwyeF-1}\nwindexuse{\nwixident{Formula.is{\_}imp}}{Formula.is:unimp}{NW1mR7zJ-4HwyeF-1}\nwindexuse{\nwixident{Formula.is{\_}not}}{Formula.is:unnot}{NW1mR7zJ-4HwyeF-1}\nwindexuse{\nwixident{Formula.is{\_}or}}{Formula.is:unor}{NW1mR7zJ-4HwyeF-1}\nwindexuse{\nwixident{Formula.mk{\_}and}}{Formula.mk:unand}{NW1mR7zJ-4HwyeF-1}\nwindexuse{\nwixident{Formula.mk{\_}exists}}{Formula.mk:unexists}{NW1mR7zJ-4HwyeF-1}\nwindexuse{\nwixident{Formula.mk{\_}forall}}{Formula.mk:unforall}{NW1mR7zJ-4HwyeF-1}\nwindexuse{\nwixident{Formula.mk{\_}not}}{Formula.mk:unnot}{NW1mR7zJ-4HwyeF-1}\nwindexuse{\nwixident{Formula.mk{\_}or}}{Formula.mk:unor}{NW1mR7zJ-4HwyeF-1}\nwendcode{}\nwbegindocs{8}\nwdocspar

\begin{node}
Skolemization eliminates existentially quantified variables by making
them functions of universally quantified formulas. The idea is that
the following are equivalent:
\begin{enumerate}
\item $\forall x\ldotp\exists y\ldotp P[x,y]$
\item $\forall x\ldotp P[x, f(x)]$ where $f$ is a ``skolemization function''.
\end{enumerate}
This is typically understood as requiring the axiom of choice (since
$f$ is choosing ``some'' constant, in some unprescribed way).
\end{node}

\begin{node}
We need to name Skolem functions to be something new and fresh. We
gather the name of all the functions currently used.
\end{node}

\nwenddocs{}\nwbegincode{9}\sublabel{NW1mR7zJ-43eePE-1}\nwmargintag{{\nwtagstyle{}\subpageref{NW1mR7zJ-43eePE-1}}}\moddef{Gather all the functions used~{\nwtagstyle{}\subpageref{NW1mR7zJ-43eePE-1}}}\endmoddef\nwstartdeflinemarkup\nwusesondefline{\\{NW1mR7zJ-4B0sXE-1}}\nwenddeflinemarkup
(* funcs : Term.t -> (string * Sort.t * int) list *)
fun funcs (Term.Var _) = []
  | funcs (Term.Fun(f,s,args)) =
      List.foldr (fn (arg, acc) =>
                   (funcs arg)@acc)
                 [(f,s,length args)]
                 args;

(* functions : Formula.t -> (string * Sort.t * int) list *)
fun functions fm =
  Formula.atom_union
    (fn p => let val (_,args) = Formula.dest_pred p
             in List.foldr (fn (arg, acc) => (funcs arg)@acc)
                           []
                           args
             end)
    fm;
\nwused{\\{NW1mR7zJ-4B0sXE-1}}\nwidentuses{\\{{\nwixident{Formula.atom{\_}union}}{Formula.atom:ununion}}\\{{\nwixident{Formula.dest{\_}pred}}{Formula.dest:unpred}}\\{{\nwixident{Formula.t}}{Formula.t}}\\{{\nwixident{Sort.t}}{Sort.t}}\\{{\nwixident{Term.t}}{Term.t}}}\nwindexuse{\nwixident{Formula.atom{\_}union}}{Formula.atom:ununion}{NW1mR7zJ-43eePE-1}\nwindexuse{\nwixident{Formula.dest{\_}pred}}{Formula.dest:unpred}{NW1mR7zJ-43eePE-1}\nwindexuse{\nwixident{Formula.t}}{Formula.t}{NW1mR7zJ-43eePE-1}\nwindexuse{\nwixident{Sort.t}}{Sort.t}{NW1mR7zJ-43eePE-1}\nwindexuse{\nwixident{Term.t}}{Term.t}{NW1mR7zJ-43eePE-1}\nwendcode{}\nwbegindocs{10}\nwdocspar

\begin{node}
We will want to take the variant of an identifier to make sure it is
not currently used for a function. We ignore the arity of the function
name, we just need a ``fresh'' function name.
\end{node}

\nwenddocs{}\nwbegincode{11}\sublabel{NW1mR7zJ-2QPys7-1}\nwmargintag{{\nwtagstyle{}\subpageref{NW1mR7zJ-2QPys7-1}}}\moddef{Variant of an identifier~{\nwtagstyle{}\subpageref{NW1mR7zJ-2QPys7-1}}}\endmoddef\nwstartdeflinemarkup\nwusesondefline{\\{NW1mR7zJ-4B0sXE-1}}\nwenddeflinemarkup
(* variant : string -> Sort.t -> (string * Sort.t * int) list -> string *)
local
  fun iter n name sort funcs =
    if List.exists (fn (x,s,_) =>
                       s = sort andalso
                       (name ^ (Int.toString n)) = x)
                   funcs
    then iter (n + 1) name sort funcs
    else (name ^ (Int.toString n), sort);
in
fun variant name sort funcs : string * Sort.t =
  if List.exists (fn (x,s,_) => s = sort andalso name = x)
                 funcs
  then iter 0 name sort funcs
  else (name, sort);
end;
\nwused{\\{NW1mR7zJ-4B0sXE-1}}\nwidentuses{\\{{\nwixident{Sort.t}}{Sort.t}}}\nwindexuse{\nwixident{Sort.t}}{Sort.t}{NW1mR7zJ-2QPys7-1}\nwendcode{}\nwbegindocs{12}\nwdocspar

\begin{node}
Skolemization acts on closed formulas in negation normal form.
So negation has been pulled down to the literals. The basic steps in the procedure is:
\begin{enumerate}
\item Given $\exists x\ldotp A[x]$, collect the distinct free
  variables in a list $(x_{1},\dots,x_{n})$ (which, since we assumed
  the given formula is closed, means the free variables are
  universally quantified and this is being recursively called on a
  subformula), then introduce the Skolem
  function $f(x_{1},\dots,x_{n})$ and return $A[f(x_{1},\dots,x_{n})]$
\item Given $A\land B$, simply return the conjunction of the
  {\Tt{}skolem\nwendquote} function applied to $A$ and $B$
\item Given $A\lor B$, simply return the disjunction of the
  {\Tt{}skolem\nwendquote} function applied to $A$ and $B$
\item Given $\forall x\ldotp A[x]$, recursively apply {\Tt{}skolem\nwendquote} to
  $A$ to obtain $A'[x]$, then return $\forall x\ldotp A'[x]$.
\end{enumerate}
\end{node}

\nwenddocs{}\nwbegincode{13}\sublabel{NW1mR7zJ-4B0sXE-1}\nwmargintag{{\nwtagstyle{}\subpageref{NW1mR7zJ-4B0sXE-1}}}\moddef{Skolemize a formula~{\nwtagstyle{}\subpageref{NW1mR7zJ-4B0sXE-1}}}\endmoddef\nwstartdeflinemarkup\nwusesondefline{\\{NW1mR7zJ-4QrsTb-1}}\nwenddeflinemarkup
\LA{}Gather all the functions used~{\nwtagstyle{}\subpageref{NW1mR7zJ-43eePE-1}}\RA{}

\LA{}Variant of an identifier~{\nwtagstyle{}\subpageref{NW1mR7zJ-2QPys7-1}}\RA{}

local
  fun iter [] acc = acc
    | iter (x::xs) acc = if List.exists (Term.eq x) xs
                         then iter xs acc
                         else iter xs (x::acc);
in
fun dedupe vars = iter vars []
end;

(* skolem : Formula.t ->  (string * Sort.t * int) list ->
     Formula.t * (string * Sort.t * int) list

REQUIRES: given formula is in NNF
ENSURES: resulting formula is Skolemized, and produces the
         list of Skolem functions/constants. *)
fun skolem (fm : Formula.t) (fns : (string * Sort.t * int) list)
  : Formula.t * ((string * Sort.t * int) list) =
  if Formula.is_exists fm
  then let val((var as (Term.Var(y,s))), A) = Formula.dest_exists fm;
           val xs = dedupe(Formula.fv fm);
           val (const as (f,s')) = variant (if List.null xs
                                 then ("c_"^y)
                                 else ("f_"^y))
                                s
                                fns;
           val fx = Term.Fun(f,s,xs);
       in skolem (Formula.subst var fx A) ((f,s',length xs)::fns) end
  else if Formula.is_forall fm
  then let val (x, A) = Formula.dest_forall fm;
           val (A', fns') = skolem A fns;
       in (Formula.mk_forall(x, A'), fns') end
  else if Formula.is_and fm
  then let val (A,B) = Formula.dest_and fm
       in skolem2 Formula.mk_and (A,B) fns end
  else if Formula.is_or fm
  then let val (A,B) = Formula.dest_or fm
       in skolem2 Formula.mk_or (A,B) fns end
  else (fm,fns)
and skolem2 conn (A,B) fns =
    let val (A',fns') = skolem A fns;
        val (B',fns'') = skolem B fns';
    in (conn(A', B'), fns'') end;
\nwused{\\{NW1mR7zJ-4QrsTb-1}}\nwidentuses{\\{{\nwixident{Formula.dest{\_}and}}{Formula.dest:unand}}\\{{\nwixident{Formula.dest{\_}or}}{Formula.dest:unor}}\\{{\nwixident{Formula.fv}}{Formula.fv}}\\{{\nwixident{Formula.is{\_}and}}{Formula.is:unand}}\\{{\nwixident{Formula.is{\_}or}}{Formula.is:unor}}\\{{\nwixident{Formula.mk{\_}and}}{Formula.mk:unand}}\\{{\nwixident{Formula.mk{\_}forall}}{Formula.mk:unforall}}\\{{\nwixident{Formula.mk{\_}or}}{Formula.mk:unor}}\\{{\nwixident{Formula.subst}}{Formula.subst}}\\{{\nwixident{Formula.t}}{Formula.t}}\\{{\nwixident{Sort.t}}{Sort.t}}\\{{\nwixident{Term.eq}}{Term.eq}}}\nwindexuse{\nwixident{Formula.dest{\_}and}}{Formula.dest:unand}{NW1mR7zJ-4B0sXE-1}\nwindexuse{\nwixident{Formula.dest{\_}or}}{Formula.dest:unor}{NW1mR7zJ-4B0sXE-1}\nwindexuse{\nwixident{Formula.fv}}{Formula.fv}{NW1mR7zJ-4B0sXE-1}\nwindexuse{\nwixident{Formula.is{\_}and}}{Formula.is:unand}{NW1mR7zJ-4B0sXE-1}\nwindexuse{\nwixident{Formula.is{\_}or}}{Formula.is:unor}{NW1mR7zJ-4B0sXE-1}\nwindexuse{\nwixident{Formula.mk{\_}and}}{Formula.mk:unand}{NW1mR7zJ-4B0sXE-1}\nwindexuse{\nwixident{Formula.mk{\_}forall}}{Formula.mk:unforall}{NW1mR7zJ-4B0sXE-1}\nwindexuse{\nwixident{Formula.mk{\_}or}}{Formula.mk:unor}{NW1mR7zJ-4B0sXE-1}\nwindexuse{\nwixident{Formula.subst}}{Formula.subst}{NW1mR7zJ-4B0sXE-1}\nwindexuse{\nwixident{Formula.t}}{Formula.t}{NW1mR7zJ-4B0sXE-1}\nwindexuse{\nwixident{Sort.t}}{Sort.t}{NW1mR7zJ-4B0sXE-1}\nwindexuse{\nwixident{Term.eq}}{Term.eq}{NW1mR7zJ-4B0sXE-1}\nwendcode{}\nwbegindocs{14}\nwdocspar

\begin{node}
We also want to simplify formulas using various identities. The basic
structure of simplification amounts to performing these
transformations once using the {\Tt{}simplify1\nwendquote} function, and propagate
them down subformulas using the {\Tt{}simplify\nwendquote} function.
\end{node}

\nwenddocs{}\nwbegincode{15}\sublabel{NW1mR7zJ-2RorhN-1}\nwmargintag{{\nwtagstyle{}\subpageref{NW1mR7zJ-2RorhN-1}}}\moddef{Simplify formulas using classical tautologies~{\nwtagstyle{}\subpageref{NW1mR7zJ-2RorhN-1}}}\endmoddef\nwstartdeflinemarkup\nwusesondefline{\\{NW1mR7zJ-l5jX3-1}}\nwenddeflinemarkup
local
  fun simplify1 fm =
    \LA{}Simplify the body of a quantified formula~{\nwtagstyle{}\subpageref{NW1mR7zJ-3pvELX-1}}\RA{}
    \LA{}Simplify double negation~{\nwtagstyle{}\subpageref{NW1mR7zJ-tCoNW-1}}\RA{}
    \LA{}Simplify conjunction~{\nwtagstyle{}\subpageref{NW1mR7zJ-19fhtN-1}}\RA{}
    \LA{}Simplify disjunction~{\nwtagstyle{}\subpageref{NW1mR7zJ-4ZXXl0-1}}\RA{}
    \LA{}Simplify implication~{\nwtagstyle{}\subpageref{NW1mR7zJ-3HyT53-1}}\RA{}
    \LA{}Simplify logical equivalence~{\nwtagstyle{}\subpageref{NW1mR7zJ-48args-1}}\RA{}
    else fm;
in
fun simplify fm =
  if Formula.is_not fm
  then let val A = Formula.dest_not fm
       in simplify1(Formula.mk_not(simplify A)) end
  else if Formula.is_and fm
  then let val (A,B) = Formula.dest_and fm
       in simplify1(Formula.mk_and(simplify A,simplify B)) end
  else if Formula.is_or fm
  then let val (A,B) = Formula.dest_or fm
       in simplify1(Formula.mk_or(simplify A,simplify B)) end
  else if Formula.is_imp fm
  then let val (A,B) = Formula.dest_imp fm
       in simplify1(Formula.mk_imp(simplify A,simplify B)) end
  else if Formula.is_iff fm
  then let val (A,B) = Formula.dest_iff fm
       in simplify1(Formula.mk_iff(simplify A,simplify B)) end
  else if Formula.is_exists fm
  then let val (x,A) = Formula.dest_exists fm
       in simplify1(Formula.mk_exists(x, simplify A)) end
  else if Formula.is_forall fm
  then let val (x,A) = Formula.dest_forall fm
       in simplify1(Formula.mk_forall(x, simplify A)) end
  else fm (* falsum or predicates are already simplified *)
end;
\nwused{\\{NW1mR7zJ-l5jX3-1}}\nwidentuses{\\{{\nwixident{Formula.dest{\_}and}}{Formula.dest:unand}}\\{{\nwixident{Formula.dest{\_}iff}}{Formula.dest:uniff}}\\{{\nwixident{Formula.dest{\_}imp}}{Formula.dest:unimp}}\\{{\nwixident{Formula.dest{\_}not}}{Formula.dest:unnot}}\\{{\nwixident{Formula.dest{\_}or}}{Formula.dest:unor}}\\{{\nwixident{Formula.is{\_}and}}{Formula.is:unand}}\\{{\nwixident{Formula.is{\_}iff}}{Formula.is:uniff}}\\{{\nwixident{Formula.is{\_}imp}}{Formula.is:unimp}}\\{{\nwixident{Formula.is{\_}not}}{Formula.is:unnot}}\\{{\nwixident{Formula.is{\_}or}}{Formula.is:unor}}\\{{\nwixident{Formula.mk{\_}and}}{Formula.mk:unand}}\\{{\nwixident{Formula.mk{\_}exists}}{Formula.mk:unexists}}\\{{\nwixident{Formula.mk{\_}forall}}{Formula.mk:unforall}}\\{{\nwixident{Formula.mk{\_}iff}}{Formula.mk:uniff}}\\{{\nwixident{Formula.mk{\_}imp}}{Formula.mk:unimp}}\\{{\nwixident{Formula.mk{\_}not}}{Formula.mk:unnot}}\\{{\nwixident{Formula.mk{\_}or}}{Formula.mk:unor}}}\nwindexuse{\nwixident{Formula.dest{\_}and}}{Formula.dest:unand}{NW1mR7zJ-2RorhN-1}\nwindexuse{\nwixident{Formula.dest{\_}iff}}{Formula.dest:uniff}{NW1mR7zJ-2RorhN-1}\nwindexuse{\nwixident{Formula.dest{\_}imp}}{Formula.dest:unimp}{NW1mR7zJ-2RorhN-1}\nwindexuse{\nwixident{Formula.dest{\_}not}}{Formula.dest:unnot}{NW1mR7zJ-2RorhN-1}\nwindexuse{\nwixident{Formula.dest{\_}or}}{Formula.dest:unor}{NW1mR7zJ-2RorhN-1}\nwindexuse{\nwixident{Formula.is{\_}and}}{Formula.is:unand}{NW1mR7zJ-2RorhN-1}\nwindexuse{\nwixident{Formula.is{\_}iff}}{Formula.is:uniff}{NW1mR7zJ-2RorhN-1}\nwindexuse{\nwixident{Formula.is{\_}imp}}{Formula.is:unimp}{NW1mR7zJ-2RorhN-1}\nwindexuse{\nwixident{Formula.is{\_}not}}{Formula.is:unnot}{NW1mR7zJ-2RorhN-1}\nwindexuse{\nwixident{Formula.is{\_}or}}{Formula.is:unor}{NW1mR7zJ-2RorhN-1}\nwindexuse{\nwixident{Formula.mk{\_}and}}{Formula.mk:unand}{NW1mR7zJ-2RorhN-1}\nwindexuse{\nwixident{Formula.mk{\_}exists}}{Formula.mk:unexists}{NW1mR7zJ-2RorhN-1}\nwindexuse{\nwixident{Formula.mk{\_}forall}}{Formula.mk:unforall}{NW1mR7zJ-2RorhN-1}\nwindexuse{\nwixident{Formula.mk{\_}iff}}{Formula.mk:uniff}{NW1mR7zJ-2RorhN-1}\nwindexuse{\nwixident{Formula.mk{\_}imp}}{Formula.mk:unimp}{NW1mR7zJ-2RorhN-1}\nwindexuse{\nwixident{Formula.mk{\_}not}}{Formula.mk:unnot}{NW1mR7zJ-2RorhN-1}\nwindexuse{\nwixident{Formula.mk{\_}or}}{Formula.mk:unor}{NW1mR7zJ-2RorhN-1}\nwendcode{}\nwbegindocs{16}\nwdocspar

\begin{node}
We simplify $\forall x\ldotp A[x]$ by simplifying $A[x]$ to $A'[x]$
(whatever that might be), then returning $\forall x\ldotp A'[x]$.
Similar simplification applies to existentially quantified formulas.
\end{node}

\nwenddocs{}\nwbegincode{17}\sublabel{NW1mR7zJ-3pvELX-1}\nwmargintag{{\nwtagstyle{}\subpageref{NW1mR7zJ-3pvELX-1}}}\moddef{Simplify the body of a quantified formula~{\nwtagstyle{}\subpageref{NW1mR7zJ-3pvELX-1}}}\endmoddef\nwstartdeflinemarkup\nwusesondefline{\\{NW1mR7zJ-2RorhN-1}}\nwenddeflinemarkup
if Formula.is_forall fm
then let val (x,A) = Formula.dest_forall fm
     in if List.exists (Term.eq x) (Formula.fv A)
        then fm
        else A
     end
else if Formula.is_exists fm
then let val (x,A) = Formula.dest_exists fm
     in if List.exists (Term.eq x) (Formula.fv A)
        then fm
        else A
     end
\nwused{\\{NW1mR7zJ-2RorhN-1}}\nwidentuses{\\{{\nwixident{Formula.fv}}{Formula.fv}}\\{{\nwixident{Term.eq}}{Term.eq}}}\nwindexuse{\nwixident{Formula.fv}}{Formula.fv}{NW1mR7zJ-3pvELX-1}\nwindexuse{\nwixident{Term.eq}}{Term.eq}{NW1mR7zJ-3pvELX-1}\nwendcode{}\nwbegindocs{18}\nwdocspar

\begin{node}
We simplify $\neg\neg A$ to be just $A$.
\end{node}

\nwenddocs{}\nwbegincode{19}\sublabel{NW1mR7zJ-tCoNW-1}\nwmargintag{{\nwtagstyle{}\subpageref{NW1mR7zJ-tCoNW-1}}}\moddef{Simplify double negation~{\nwtagstyle{}\subpageref{NW1mR7zJ-tCoNW-1}}}\endmoddef\nwstartdeflinemarkup\nwusesondefline{\\{NW1mR7zJ-2RorhN-1}}\nwenddeflinemarkup
else if Formula.is_not fm
then let val A = Formula.dest_not fm
     in if Formula.is_not A
        then Formula.dest_not A
        else fm
     end
\nwused{\\{NW1mR7zJ-2RorhN-1}}\nwidentuses{\\{{\nwixident{Formula.dest{\_}not}}{Formula.dest:unnot}}\\{{\nwixident{Formula.is{\_}not}}{Formula.is:unnot}}}\nwindexuse{\nwixident{Formula.dest{\_}not}}{Formula.dest:unnot}{NW1mR7zJ-tCoNW-1}\nwindexuse{\nwixident{Formula.is{\_}not}}{Formula.is:unnot}{NW1mR7zJ-tCoNW-1}\nwendcode{}\nwbegindocs{20}\nwdocspar

\begin{node}
We simplify $\bot\land B$ and $A\land\bot$ to $\bot$. In the cases
where neither conjunct is falsum: 
$\top\land B$ (i.e., $A=\top$) simplifies to $B$, and $A\land\top$
(i.e., $B=\top$) simplifies to $A$.
\end{node}

\nwenddocs{}\nwbegincode{21}\sublabel{NW1mR7zJ-19fhtN-1}\nwmargintag{{\nwtagstyle{}\subpageref{NW1mR7zJ-19fhtN-1}}}\moddef{Simplify conjunction~{\nwtagstyle{}\subpageref{NW1mR7zJ-19fhtN-1}}}\endmoddef\nwstartdeflinemarkup\nwusesondefline{\\{NW1mR7zJ-2RorhN-1}}\nwenddeflinemarkup
else if Formula.is_and fm
then let val (A,B) = Formula.dest_and fm
     in if Formula.is_falsum A orelse Formula.is_falsum B
        then Formula.mk_falsum ()
        else if Formula.is_verum A
        then B
        else if Formula.is_verum B
        then A
        else fm
     end
\nwused{\\{NW1mR7zJ-2RorhN-1}}\nwidentuses{\\{{\nwixident{Formula.dest{\_}and}}{Formula.dest:unand}}\\{{\nwixident{Formula.is{\_}and}}{Formula.is:unand}}\\{{\nwixident{Formula.is{\_}falsum}}{Formula.is:unfalsum}}\\{{\nwixident{Formula.is{\_}verum}}{Formula.is:unverum}}\\{{\nwixident{Formula.mk{\_}falsum}}{Formula.mk:unfalsum}}}\nwindexuse{\nwixident{Formula.dest{\_}and}}{Formula.dest:unand}{NW1mR7zJ-19fhtN-1}\nwindexuse{\nwixident{Formula.is{\_}and}}{Formula.is:unand}{NW1mR7zJ-19fhtN-1}\nwindexuse{\nwixident{Formula.is{\_}falsum}}{Formula.is:unfalsum}{NW1mR7zJ-19fhtN-1}\nwindexuse{\nwixident{Formula.is{\_}verum}}{Formula.is:unverum}{NW1mR7zJ-19fhtN-1}\nwindexuse{\nwixident{Formula.mk{\_}falsum}}{Formula.mk:unfalsum}{NW1mR7zJ-19fhtN-1}\nwendcode{}\nwbegindocs{22}\nwdocspar

\begin{node}
We simplify $\bot\lor A$ and $A\lor\bot$ to $A$, and $\top\lor A$
and $A\lor\top$ to $\top$. Otherwise we don't touch the formula.
\end{node}

\nwenddocs{}\nwbegincode{23}\sublabel{NW1mR7zJ-4ZXXl0-1}\nwmargintag{{\nwtagstyle{}\subpageref{NW1mR7zJ-4ZXXl0-1}}}\moddef{Simplify disjunction~{\nwtagstyle{}\subpageref{NW1mR7zJ-4ZXXl0-1}}}\endmoddef\nwstartdeflinemarkup\nwusesondefline{\\{NW1mR7zJ-2RorhN-1}}\nwenddeflinemarkup
else if Formula.is_or fm
then let val (A,B) = Formula.dest_or fm
     in if Formula.is_falsum A orelse Formula.is_verum B
        then B
        else if Formula.is_falsum B orelse Formula.is_verum A
        then A
        else fm
     end
\nwused{\\{NW1mR7zJ-2RorhN-1}}\nwidentuses{\\{{\nwixident{Formula.dest{\_}or}}{Formula.dest:unor}}\\{{\nwixident{Formula.is{\_}falsum}}{Formula.is:unfalsum}}\\{{\nwixident{Formula.is{\_}or}}{Formula.is:unor}}\\{{\nwixident{Formula.is{\_}verum}}{Formula.is:unverum}}}\nwindexuse{\nwixident{Formula.dest{\_}or}}{Formula.dest:unor}{NW1mR7zJ-4ZXXl0-1}\nwindexuse{\nwixident{Formula.is{\_}falsum}}{Formula.is:unfalsum}{NW1mR7zJ-4ZXXl0-1}\nwindexuse{\nwixident{Formula.is{\_}or}}{Formula.is:unor}{NW1mR7zJ-4ZXXl0-1}\nwindexuse{\nwixident{Formula.is{\_}verum}}{Formula.is:unverum}{NW1mR7zJ-4ZXXl0-1}\nwendcode{}\nwbegindocs{24}\nwdocspar

\begin{node}
We simplify $\bot\implies A$ and $A\implies\top$ both to just $\top$.
We similarly simplify $\top\implies B$ to just $B$. Otherwise we leave
the formula alone.
\end{node}

\nwenddocs{}\nwbegincode{25}\sublabel{NW1mR7zJ-3HyT53-1}\nwmargintag{{\nwtagstyle{}\subpageref{NW1mR7zJ-3HyT53-1}}}\moddef{Simplify implication~{\nwtagstyle{}\subpageref{NW1mR7zJ-3HyT53-1}}}\endmoddef\nwstartdeflinemarkup\nwusesondefline{\\{NW1mR7zJ-2RorhN-1}}\nwenddeflinemarkup
else if Formula.is_imp fm
then let val (A,B) = Formula.dest_imp fm
     in if Formula.is_falsum A orelse Formula.is_verum B
        then Formula.mk_verum ()
        else if Formula.is_verum A
        then B
        else fm
     end
\nwused{\\{NW1mR7zJ-2RorhN-1}}\nwidentuses{\\{{\nwixident{Formula.dest{\_}imp}}{Formula.dest:unimp}}\\{{\nwixident{Formula.is{\_}falsum}}{Formula.is:unfalsum}}\\{{\nwixident{Formula.is{\_}imp}}{Formula.is:unimp}}\\{{\nwixident{Formula.is{\_}verum}}{Formula.is:unverum}}\\{{\nwixident{Formula.mk{\_}verum}}{Formula.mk:unverum}}}\nwindexuse{\nwixident{Formula.dest{\_}imp}}{Formula.dest:unimp}{NW1mR7zJ-3HyT53-1}\nwindexuse{\nwixident{Formula.is{\_}falsum}}{Formula.is:unfalsum}{NW1mR7zJ-3HyT53-1}\nwindexuse{\nwixident{Formula.is{\_}imp}}{Formula.is:unimp}{NW1mR7zJ-3HyT53-1}\nwindexuse{\nwixident{Formula.is{\_}verum}}{Formula.is:unverum}{NW1mR7zJ-3HyT53-1}\nwindexuse{\nwixident{Formula.mk{\_}verum}}{Formula.mk:unverum}{NW1mR7zJ-3HyT53-1}\nwendcode{}\nwbegindocs{26}\nwdocspar

\begin{node}
We simplify $A\iff B$ according to the following rules:
\begin{enumerate}
\item $\top\iff B$ becomes $B$
\item $A\iff\top$ becomes $A$
\item $\bot\iff B$ becomes $\neg B$
\item $A\iff\bot$ becomes $\neg A$
\end{enumerate}
Otherwise, we just leave the formula alone.
\end{node}

\nwenddocs{}\nwbegincode{27}\sublabel{NW1mR7zJ-48args-1}\nwmargintag{{\nwtagstyle{}\subpageref{NW1mR7zJ-48args-1}}}\moddef{Simplify logical equivalence~{\nwtagstyle{}\subpageref{NW1mR7zJ-48args-1}}}\endmoddef\nwstartdeflinemarkup\nwusesondefline{\\{NW1mR7zJ-2RorhN-1}}\nwenddeflinemarkup
else if Formula.is_iff fm
then let val (A,B) = Formula.dest_iff fm
     in if Formula.is_verum A
        then B
        else if Formula.is_verum B
        then A
        else if Formula.is_falsum A
        then Formula.mk_not B
        else if Formula.is_falsum B
        then Formula.mk_not A
        else fm
     end
\nwused{\\{NW1mR7zJ-2RorhN-1}}\nwidentuses{\\{{\nwixident{Formula.dest{\_}iff}}{Formula.dest:uniff}}\\{{\nwixident{Formula.is{\_}falsum}}{Formula.is:unfalsum}}\\{{\nwixident{Formula.is{\_}iff}}{Formula.is:uniff}}\\{{\nwixident{Formula.is{\_}verum}}{Formula.is:unverum}}\\{{\nwixident{Formula.mk{\_}not}}{Formula.mk:unnot}}}\nwindexuse{\nwixident{Formula.dest{\_}iff}}{Formula.dest:uniff}{NW1mR7zJ-48args-1}\nwindexuse{\nwixident{Formula.is{\_}falsum}}{Formula.is:unfalsum}{NW1mR7zJ-48args-1}\nwindexuse{\nwixident{Formula.is{\_}iff}}{Formula.is:uniff}{NW1mR7zJ-48args-1}\nwindexuse{\nwixident{Formula.is{\_}verum}}{Formula.is:unverum}{NW1mR7zJ-48args-1}\nwindexuse{\nwixident{Formula.mk{\_}not}}{Formula.mk:unnot}{NW1mR7zJ-48args-1}\nwendcode{}\nwbegindocs{28}\nwdocspar

\begin{node}
Skolemization then applies the algorithm to replace existentially
quantified variables with skolem constants, to the negation normal
form of the given formula (after simplifying the formula).
\end{node}

\nwenddocs{}\nwbegincode{29}\sublabel{NW1mR7zJ-4QrsTb-1}\nwmargintag{{\nwtagstyle{}\subpageref{NW1mR7zJ-4QrsTb-1}}}\moddef{Skolemization~{\nwtagstyle{}\subpageref{NW1mR7zJ-4QrsTb-1}}}\endmoddef\nwstartdeflinemarkup\nwusesondefline{\\{NW1mR7zJ-l5jX3-1}}\nwenddeflinemarkup
\LA{}Skolemize a formula~{\nwtagstyle{}\subpageref{NW1mR7zJ-4B0sXE-1}}\RA{}

(* askolemize : Formula.t -> Formula.t

REQUIRES: The given formula is a closed formula.
ENSURES: The resulting formula is in Negation Normal Form with
         all existential quantifiers replaced by Skolem
         functions/constants. *)
fun askolemize fm =
  let val (result, _) = skolem (nnf(simplify(fm))) (functions fm)
  in result end;
\nwused{\\{NW1mR7zJ-l5jX3-1}}\nwidentuses{\\{{\nwixident{Formula.t}}{Formula.t}}}\nwindexuse{\nwixident{Formula.t}}{Formula.t}{NW1mR7zJ-4QrsTb-1}\nwendcode{}\nwbegindocs{30}\nwdocspar

\section{Tableaux}

\begin{node}
Most tableaux methods follows \leanTAP's algorithm~\cite{beckert1995leantap}.
The basic tableaux method works as follows:
\begin{enumerate}
\item for $A\land B$, separately assume $A$ and $B$; this amounts to
  working through $A$, then working through $B$, then working through
  the remaining unexplored subformulas.
\item for $A\lor B$, we ``branch'' --- that is to say, we perform two
  separate refutations, one for $A$, and another for $B$;
\item for $\forall x\ldotp A[x]$, introduce a fresh variable $y$ and
  assume $A[y]$, but we keep the original $\forall x\ldotp A[x]$
  in case multiple instances are needed.
\item for a literal $L$, try to refute it based on the literals
  already encountered (which is the only way to close the branch);
  otherwise, we add $L$ to the ``database'' of
  literals, and look at the next unexpanded formula.
\end{enumerate}
There is no ordering to the subformulas being explored. Smullyan noted
that it's probably more optimal to explore the conjunctions first,
then the disjunctions. This would probably be low hanging fruit.
\end{node}

\nwenddocs{}\nwbegincode{31}\sublabel{NW1mR7zJ-1k5nOE-1}\nwmargintag{{\nwtagstyle{}\subpageref{NW1mR7zJ-1k5nOE-1}}}\moddef{Tableau method~{\nwtagstyle{}\subpageref{NW1mR7zJ-1k5nOE-1}}}\endmoddef\nwstartdeflinemarkup\nwusesondefline{\\{NW1mR7zJ-yCIr5-1}}\nwenddeflinemarkup
fun tryfind f [] = raise Fail "tryfind"
  | tryfind f (x::xs) = (f x
                         handle _ => tryfind f xs);

(* tableau : (Formula.t list * Formula.t list * int) ->
                (Unif.Env.t * int -> 'a) ->
                Unif.Env.t * int -> 'a *)
fun tableau (fms, lits, n) (cont : Unif.Env.t * int -> 'a) (env, k) =
  if n < 0 then raise Fail "no proof at this level"
  else
    (case fms of
      [] => raise Fail "tableau: no proof"
    | (phi::unexp) =>
        if Formula.is_and phi
        then let val (A,B) = Formula.dest_and phi
             in tableau (A::B::unexp,lits,n) cont (env, k) end
        else if Formula.is_or phi
        then let val (A,B) = Formula.dest_or phi
             in tableau (A::unexp,lits,n)
                        (tableau (B::unexp,lits,n) cont)
                        (env, k)
             end
        else if Formula.is_forall phi
        then let val(x,A) = Formula.dest_forall phi;
                 val y = Term.Var("_"^(Int.toString k), Term.sort(x));
                 val A' = Formula.subst x y A;
             in tableau (A'::unexp@[Formula.mk_forall(x,A)],lits,n-1) cont (env, k+1) end
        else
         (tryfind (fn l => cont(Unif.y_complements env (phi,l),k))
                  lits
         handle _ => tableau(unexp,phi::lits,n) cont (env,k)));
\nwused{\\{NW1mR7zJ-yCIr5-1}}\nwidentuses{\\{{\nwixident{Env.t}}{Env.t}}\\{{\nwixident{Formula.dest{\_}and}}{Formula.dest:unand}}\\{{\nwixident{Formula.dest{\_}or}}{Formula.dest:unor}}\\{{\nwixident{Formula.is{\_}and}}{Formula.is:unand}}\\{{\nwixident{Formula.is{\_}or}}{Formula.is:unor}}\\{{\nwixident{Formula.mk{\_}forall}}{Formula.mk:unforall}}\\{{\nwixident{Formula.subst}}{Formula.subst}}\\{{\nwixident{Formula.t}}{Formula.t}}\\{{\nwixident{Term.sort}}{Term.sort}}\\{{\nwixident{Unif.y{\_}complements}}{Unif.y:uncomplements}}}\nwindexuse{\nwixident{Env.t}}{Env.t}{NW1mR7zJ-1k5nOE-1}\nwindexuse{\nwixident{Formula.dest{\_}and}}{Formula.dest:unand}{NW1mR7zJ-1k5nOE-1}\nwindexuse{\nwixident{Formula.dest{\_}or}}{Formula.dest:unor}{NW1mR7zJ-1k5nOE-1}\nwindexuse{\nwixident{Formula.is{\_}and}}{Formula.is:unand}{NW1mR7zJ-1k5nOE-1}\nwindexuse{\nwixident{Formula.is{\_}or}}{Formula.is:unor}{NW1mR7zJ-1k5nOE-1}\nwindexuse{\nwixident{Formula.mk{\_}forall}}{Formula.mk:unforall}{NW1mR7zJ-1k5nOE-1}\nwindexuse{\nwixident{Formula.subst}}{Formula.subst}{NW1mR7zJ-1k5nOE-1}\nwindexuse{\nwixident{Formula.t}}{Formula.t}{NW1mR7zJ-1k5nOE-1}\nwindexuse{\nwixident{Term.sort}}{Term.sort}{NW1mR7zJ-1k5nOE-1}\nwindexuse{\nwixident{Unif.y{\_}complements}}{Unif.y:uncomplements}{NW1mR7zJ-1k5nOE-1}\nwendcode{}\nwbegindocs{32}\nwdocspar

\begin{node}
The basic idea is to proving a formula $\phi$ classically with method
of tableaux amounts to refuting $\neg\phi$ with tableaux. We need to
universally quantify free variables, then negate the result, then
Skolemize the result, before applying the tableaux method.
\end{node}

\nwenddocs{}\nwbegincode{33}\sublabel{NW1mR7zJ-3SEP5w-1}\nwmargintag{{\nwtagstyle{}\subpageref{NW1mR7zJ-3SEP5w-1}}}\moddef{Tableaux refutation~{\nwtagstyle{}\subpageref{NW1mR7zJ-3SEP5w-1}}}\endmoddef\nwstartdeflinemarkup\nwusesondefline{\\{NW1mR7zJ-yCIr5-1}}\nwenddeflinemarkup
fun deepen f n maxiter =
  (print("Searching with depth limit "^
         (Int.toString n)^"\\n");
   f n)
  handle (Fail _) => if n < maxiter then deepen f (n + 1) maxiter
                     else raise Fail "maximum iterations encountered";

local
  fun id x = x;
in
fun tabrefute fms maxiter =
  deepen (fn n =>
             (tableau (fms,[],n)
                      id
                      (Unif.Env.undefined (),0);
              n))
         0
         maxiter
end;
\nwused{\\{NW1mR7zJ-yCIr5-1}}\nwidentuses{\\{{\nwixident{Env.undefined}}{Env.undefined}}}\nwindexuse{\nwixident{Env.undefined}}{Env.undefined}{NW1mR7zJ-3SEP5w-1}\nwendcode{}\nwbegindocs{34}\nwdocspar

\begin{node}
We can combine everything to give the user {\Tt{}\nwlinkedidentq{Classical.tab}{NW1mR7zJ-yCIr5-1}\nwendquote}, which
attempts to prove the validity of the given formula by refuting its
negation. The integer parameter supplied limits the number of
iterations performed, to guarantee it will terminate.

One possibly optimization is to transform the Skolemized formula into
Disjunctive normal form, then applying the tableaux method to each
disjunct.
\end{node}

\nwenddocs{}\nwbegincode{35}\sublabel{NW1mR7zJ-yCIr5-1}\nwmargintag{{\nwtagstyle{}\subpageref{NW1mR7zJ-yCIr5-1}}}\moddef{Tableaux refuter~{\nwtagstyle{}\subpageref{NW1mR7zJ-yCIr5-1}}}\endmoddef\nwstartdeflinemarkup\nwusesondefline{\\{NW1mR7zJ-l5jX3-1}}\nwenddeflinemarkup
\LA{}Tableau method~{\nwtagstyle{}\subpageref{NW1mR7zJ-1k5nOE-1}}\RA{}

\LA{}Tableaux refutation~{\nwtagstyle{}\subpageref{NW1mR7zJ-3SEP5w-1}}\RA{}

(* generalize : Formula.t -> Formula.t *)
fun generalize fm =
  List.foldr Formula.mk_forall fm (dedupe (Formula.fv fm));

fun tab fm maxiter =
  let val fm' = askolemize(Formula.mk_not(generalize fm))
  in if Formula.is_falsum fm' then 0
     else tabrefute [fm'] maxiter
  end;
\nwindexdefn{\nwixident{Classical.tab}}{Classical.tab}{NW1mR7zJ-yCIr5-1}\eatline
\nwused{\\{NW1mR7zJ-l5jX3-1}}\nwidentdefs{\\{{\nwixident{Classical.tab}}{Classical.tab}}}\nwidentuses{\\{{\nwixident{Formula.fv}}{Formula.fv}}\\{{\nwixident{Formula.is{\_}falsum}}{Formula.is:unfalsum}}\\{{\nwixident{Formula.mk{\_}forall}}{Formula.mk:unforall}}\\{{\nwixident{Formula.mk{\_}not}}{Formula.mk:unnot}}\\{{\nwixident{Formula.t}}{Formula.t}}}\nwindexuse{\nwixident{Formula.fv}}{Formula.fv}{NW1mR7zJ-yCIr5-1}\nwindexuse{\nwixident{Formula.is{\_}falsum}}{Formula.is:unfalsum}{NW1mR7zJ-yCIr5-1}\nwindexuse{\nwixident{Formula.mk{\_}forall}}{Formula.mk:unforall}{NW1mR7zJ-yCIr5-1}\nwindexuse{\nwixident{Formula.mk{\_}not}}{Formula.mk:unnot}{NW1mR7zJ-yCIr5-1}\nwindexuse{\nwixident{Formula.t}}{Formula.t}{NW1mR7zJ-yCIr5-1}\nwendcode{}\nwbegindocs{36}\nwdocspar
\begin{node}
As a couple of unit tests, we should check that the simpler problems
found in Pelletier's problem set ``work''.
\end{node}

\nwenddocs{}\nwbegincode{37}\sublabel{NW1mR7zJ-4QAFgj-1}\nwmargintag{{\nwtagstyle{}\subpageref{NW1mR7zJ-4QAFgj-1}}}\moddef{Unit tests for \code{}Classical\edoc{}~{\nwtagstyle{}\subpageref{NW1mR7zJ-4QAFgj-1}}}\endmoddef\nwstartdeflinemarkup\nwenddeflinemarkup
local
  val A = Formula.mk_pred("A", []);
  val B = Formula.mk_pred("B", []);
  infixr ==> <=>;
  fun (p ==> q) = Formula.mk_imp(p, q);
  fun (p <=> q) = Formula.mk_iff(p, q);
  infixr &;
  fun (p & q) = Formula.mk_and(p, q);

  fun ~ p = Formula.mk_not(p);

  fun exists(x, p) = Formula.mk_exists(x, p);
  fun all(x, p) = Formula.mk_forall(x, p);
in
val p1 = \nwlinkedidentc{Classical.tab}{NW1mR7zJ-yCIr5-1} ((A ==> B) <=> ((~ B) ==> (~ A))) 1;
val p2 = \nwlinkedidentc{Classical.tab}{NW1mR7zJ-yCIr5-1} ((~(~ A)) <=> A) 1;
val p3 = \nwlinkedidentc{Classical.tab}{NW1mR7zJ-yCIr5-1} ((~(A ==> B)) ==> (A ==> B)) 1;
val p4 = \nwlinkedidentc{Classical.tab}{NW1mR7zJ-yCIr5-1} (((~ A) ==> B) <=> ((~ B) ==> A)) 1;
val p5 = \nwlinkedidentc{Classical.tab}{NW1mR7zJ-yCIr5-1} (Formula.mk_or(A, ~ A)) 1;

(* Expected: 2 iterations *)
val p18 = let val x = Term.Var("x",Sort.IND);
    val y = Term.Var("y",Sort.IND);
    fun P t = Formula.mk_pred("P", [t]);
in \nwlinkedidentc{Classical.tab}{NW1mR7zJ-yCIr5-1} (exists(y, all(x, (P y) ==> (P x)))) 3 end;

(* expected: 2 iteration *)
val p19 =
  let val x = Term.Var("x",Sort.IND);
      val y = Term.Var("y",Sort.IND);
      val z = Term.Var("z",Sort.IND);
      fun P t = Formula.mk_pred("P", [t]);
      fun Q t = Formula.mk_pred("Q", [t]);
  in \nwlinkedidentc{Classical.tab}{NW1mR7zJ-yCIr5-1} (exists(x,
                           all(y,
                               all(z, ((P y) ==> (Q z)) ==> (P x) ==> (Q x)))))
                   3
  end;

end;
\nwnotused{Unit tests for [[Classical]]}\nwidentuses{\\{{\nwixident{Classical.tab}}{Classical.tab}}\\{{\nwixident{Formula.mk{\_}and}}{Formula.mk:unand}}\\{{\nwixident{Formula.mk{\_}exists}}{Formula.mk:unexists}}\\{{\nwixident{Formula.mk{\_}forall}}{Formula.mk:unforall}}\\{{\nwixident{Formula.mk{\_}iff}}{Formula.mk:uniff}}\\{{\nwixident{Formula.mk{\_}imp}}{Formula.mk:unimp}}\\{{\nwixident{Formula.mk{\_}not}}{Formula.mk:unnot}}\\{{\nwixident{Formula.mk{\_}or}}{Formula.mk:unor}}\\{{\nwixident{Formula.mk{\_}pred}}{Formula.mk:unpred}}\\{{\nwixident{Sort.IND}}{Sort.IND}}}\nwindexuse{\nwixident{Classical.tab}}{Classical.tab}{NW1mR7zJ-4QAFgj-1}\nwindexuse{\nwixident{Formula.mk{\_}and}}{Formula.mk:unand}{NW1mR7zJ-4QAFgj-1}\nwindexuse{\nwixident{Formula.mk{\_}exists}}{Formula.mk:unexists}{NW1mR7zJ-4QAFgj-1}\nwindexuse{\nwixident{Formula.mk{\_}forall}}{Formula.mk:unforall}{NW1mR7zJ-4QAFgj-1}\nwindexuse{\nwixident{Formula.mk{\_}iff}}{Formula.mk:uniff}{NW1mR7zJ-4QAFgj-1}\nwindexuse{\nwixident{Formula.mk{\_}imp}}{Formula.mk:unimp}{NW1mR7zJ-4QAFgj-1}\nwindexuse{\nwixident{Formula.mk{\_}not}}{Formula.mk:unnot}{NW1mR7zJ-4QAFgj-1}\nwindexuse{\nwixident{Formula.mk{\_}or}}{Formula.mk:unor}{NW1mR7zJ-4QAFgj-1}\nwindexuse{\nwixident{Formula.mk{\_}pred}}{Formula.mk:unpred}{NW1mR7zJ-4QAFgj-1}\nwindexuse{\nwixident{Sort.IND}}{Sort.IND}{NW1mR7zJ-4QAFgj-1}\nwendcode{}\nwbegindocs{38}\nwdocspar

\begin{node}
I am curious whether the claim ``$\forall x\ldotp x=0\lor\exists
y,z\ldotp x=(y,z)$'' is provable in $FS_{0}$. The relevant axioms
conjoined together as the premises for the claim:
\begin{multline*}
(\forall x_{1},x_{2}\ldotp 0\neq(x_{1},x_{2}))\land
(\forall x_{3},x_{4}\ldotp P_{1}(x_{3},x_{4})=x_{3})\land
(\forall x_{4},x_{6}\ldotp P_{2}(x_{5},x_{6})=x_{6})\land\\
(0\in U)\land(\forall x_{7},x_{8}\ldotp x_{7}\in U\land x_{8}\in U\implies (x_{7},x_{8})\in U)\land\\
(\forall X\ldotp 0\in X\land\forall x_{9},x_{10}\ldotp(x_{9}\in X\land x_{10}\in X\implies(x_{9},x_{10})\in X)\implies\forall x_{11}\ldotp x_{11}\in X)\\
(\forall t_{1},t_{2}\ldotp t_{1}=t_{2}\implies t_{2}=t_{1})\land(\forall t_{3},t_{4},t_{5}\ldotp t_{3}=t_{4}\land t_{4}=t_{5}\implies t_{3}=t_{5})\\
(\forall t_{6},t_{7},f\ldotp t_{6}=t_{7}\implies ft_{6}=ft_{7})\implies\\
  \implies(\forall x_{12}\ldotp 0=x_{12}\lor\exists x_{13},x_{14}\ldotp x_{12}=(x_{13},x_{14}))
\end{multline*}
The Skolemized version of this formula is
\begin{align*}
(\forall x_{1},x_{2}\ldotp 0\neq(x_{1},x_{2}))\land\\
(\forall x_{3},x_{4}\ldotp P_{1}(x_{3},x_{4})=x_{3})\land\\
(\forall x_{4},x_{6}\ldotp P_{2}(x_{5},x_{6})=x_{6})\land\\
(0\in U)\land\\
  (\forall x_{7},x_{8}\ldotp x_{7}\notin U\lor x_{8}\notin U\lor(x_{7},x_{8})\in U)\land\\
(0\notin X\lor\bigl(f_{x}(X)\in X\land f_{y}(X)\in X\land(f_{x}(X),f_{y}(X))\notin X\bigr)\lor(\forall x\ldotp x\in X))\land\\
  (\forall t_{1},t_{2}\ldotp t_{1}\neq t_{2}\lor t_{2}=t_{1})\land\\
  (\forall t_{3},t_{4},t_{5}\ldotp t_{3}\neq t_{4}\lor t_{4}\neq t_{5}\lor t_{3}=t_{5})\land\\
(\forall t_{6},t_{7},f\ldotp t_{6}\neq t_{7}\lor ft_{6}=ft_{7})\land\\
  \land(0\neq c_{x_{12}}\land\forall x_{13},x_{14}\ldotp c_{x_{12}}\neq(x_{13},x_{14}))
\end{align*}



%% And(
%%   And(Forall(Var(x1, IND), Forall(Var(x2, IND), Not(Pred(=[Fun(cons, IND, [Var(x1, IND), Var(x2, IND)]), Fun(0, IND, [])])))),
%%   And(Forall(Var(x3, IND), Forall(Var(x4, IND), Pred(=[Fun(apply, IND, [Fun(P1, FUN, []), Fun(cons, IND, [Var(x3, IND), Var(x4, IND)])]), Var(x3, IND)]))),
%%   And(Forall(Var(x5, IND), Forall(Var(x6, IND), Pred(=[Fun(apply, IND, [Fun(P2, FUN, []), Fun(cons, IND, [Var(x5, IND), Var(x6, IND)])]), Var(x6, IND)]))),
%%   And(Pred(in[Fun(0, IND, []), Fun(U, CLASS, [])]),
%%   And(Forall(Var(x7, IND), Forall(Var(x8, IND), Or(Or(Not(Pred(in[Var(x7, IND), Fun(U, CLASS, [])])), Not(Pred(in[Var(x8, IND), Fun(U, CLASS, [])]))), Pred(in[Fun(cons, IND, [Var(x7, IND), Var(x8, IND)]), Fun(U, CLASS, [])])))),
%%   And(Forall(Var(X, CLASS), Or(Or(Not(Pred(in[Fun(0, IND, []), Var(X, CLASS)])), And(And(Pred(in[Fun(f_x, IND, [Var(X, CLASS)]), Var(X, CLASS)]), Pred(in[Fun(f_y, IND, [Var(X, CLASS)]), Var(X, CLASS)])), Not(Pred(in[Fun(cons, IND, [Fun(f_x, IND, [Var(X, CLASS)]), Fun(f_y, IND, [Var(X, CLASS)])]), Var(X, CLASS)])))), Forall(Var(x, IND), Pred(in[Var(x, IND), Var(X, CLASS)])))),
%%   And(Forall(Var(x9, IND), Forall(Var(x10, IND), Or(Not(Pred(=[Var(x9, IND), Var(x10, IND)])), Pred(=[Var(x10, IND), Var(x9, IND)])))),
%%   And(Forall(Var(x12, IND), Forall(Var(x13, IND), Forall(Var(x11, IND), Or(Or(Not(Pred(=[Var(x11, IND), Var(x12, IND)])), Not(Pred(=[Var(x12, IND), Var(x13, IND)]))), Pred(=[Var(x11, IND), Var(x13, IND)]))))),
%%       Forall(Var(x15, IND), Forall(Var(f, FUN), Forall(Var(x14, IND), Or(Not(Pred(=[Var(x14, IND), Var(x15, IND)])), Pred(=[Fun(apply, IND, [Var(f, FUN), Var(x14, IND)]), Fun(apply, IND, [Var(f, FUN), Var(x15, IND)])]))))))))))))), 
%% And(Not(Pred(=[Fun(0, IND, []), Fun(c_x16, IND, [])])),
%%     Forall(Var(x17, IND), Forall(Var(x18, IND), Not(Pred(=[Fun(c_x16, IND, []), Fun(cons, IND, [Var(x17, IND), Var(x18, IND)])]))))))
\end{node}
\nwenddocs{}

\nwixlogsorted{c}{{Classical.sig}{NW1mR7zJ-3NMRe1-1}{\nwixd{NW1mR7zJ-3NMRe1-1}}}%
\nwixlogsorted{c}{{Classical.sml}{NW1mR7zJ-l5jX3-1}{\nwixd{NW1mR7zJ-l5jX3-1}}}%
\nwixlogsorted{c}{{Gather all the functions used}{NW1mR7zJ-43eePE-1}{\nwixd{NW1mR7zJ-43eePE-1}\nwixu{NW1mR7zJ-4B0sXE-1}}}%
\nwixlogsorted{c}{{Negation normal form for negated formulas}{NW1mR7zJ-4HwyeF-1}{\nwixu{NW1mR7zJ-4TkWUc-1}\nwixd{NW1mR7zJ-4HwyeF-1}}}%
\nwixlogsorted{c}{{Simplify conjunction}{NW1mR7zJ-19fhtN-1}{\nwixu{NW1mR7zJ-2RorhN-1}\nwixd{NW1mR7zJ-19fhtN-1}}}%
\nwixlogsorted{c}{{Simplify disjunction}{NW1mR7zJ-4ZXXl0-1}{\nwixu{NW1mR7zJ-2RorhN-1}\nwixd{NW1mR7zJ-4ZXXl0-1}}}%
\nwixlogsorted{c}{{Simplify double negation}{NW1mR7zJ-tCoNW-1}{\nwixu{NW1mR7zJ-2RorhN-1}\nwixd{NW1mR7zJ-tCoNW-1}}}%
\nwixlogsorted{c}{{Simplify formulas using classical tautologies}{NW1mR7zJ-2RorhN-1}{\nwixu{NW1mR7zJ-l5jX3-1}\nwixd{NW1mR7zJ-2RorhN-1}}}%
\nwixlogsorted{c}{{Simplify implication}{NW1mR7zJ-3HyT53-1}{\nwixu{NW1mR7zJ-2RorhN-1}\nwixd{NW1mR7zJ-3HyT53-1}}}%
\nwixlogsorted{c}{{Simplify logical equivalence}{NW1mR7zJ-48args-1}{\nwixu{NW1mR7zJ-2RorhN-1}\nwixd{NW1mR7zJ-48args-1}}}%
\nwixlogsorted{c}{{Simplify the body of a quantified formula}{NW1mR7zJ-3pvELX-1}{\nwixu{NW1mR7zJ-2RorhN-1}\nwixd{NW1mR7zJ-3pvELX-1}}}%
\nwixlogsorted{c}{{Skolemization}{NW1mR7zJ-4QrsTb-1}{\nwixu{NW1mR7zJ-l5jX3-1}\nwixd{NW1mR7zJ-4QrsTb-1}}}%
\nwixlogsorted{c}{{Skolemize a formula}{NW1mR7zJ-4B0sXE-1}{\nwixd{NW1mR7zJ-4B0sXE-1}\nwixu{NW1mR7zJ-4QrsTb-1}}}%
\nwixlogsorted{c}{{Tableau method}{NW1mR7zJ-1k5nOE-1}{\nwixd{NW1mR7zJ-1k5nOE-1}\nwixu{NW1mR7zJ-yCIr5-1}}}%
\nwixlogsorted{c}{{Tableaux refutation}{NW1mR7zJ-3SEP5w-1}{\nwixd{NW1mR7zJ-3SEP5w-1}\nwixu{NW1mR7zJ-yCIr5-1}}}%
\nwixlogsorted{c}{{Tableaux refuter}{NW1mR7zJ-yCIr5-1}{\nwixu{NW1mR7zJ-l5jX3-1}\nwixd{NW1mR7zJ-yCIr5-1}}}%
\nwixlogsorted{c}{{Transform formula to negation normal form}{NW1mR7zJ-4TkWUc-1}{\nwixu{NW1mR7zJ-l5jX3-1}\nwixd{NW1mR7zJ-4TkWUc-1}}}%
\nwixlogsorted{c}{{Unit tests for \code{}Classical\edoc{}}{NW1mR7zJ-4QAFgj-1}{\nwixd{NW1mR7zJ-4QAFgj-1}}}%
\nwixlogsorted{c}{{Variant of an identifier}{NW1mR7zJ-2QPys7-1}{\nwixd{NW1mR7zJ-2QPys7-1}\nwixu{NW1mR7zJ-4B0sXE-1}}}%
\nwixlogsorted{i}{{\nwixident{Classical.tab}}{Classical.tab}}%
\nwixlogsorted{i}{{\nwixident{Env.t}}{Env.t}}%
\nwixlogsorted{i}{{\nwixident{Env.undefined}}{Env.undefined}}%
\nwixlogsorted{i}{{\nwixident{Formula.atom{\_}union}}{Formula.atom:ununion}}%
\nwixlogsorted{i}{{\nwixident{Formula.dest{\_}and}}{Formula.dest:unand}}%
\nwixlogsorted{i}{{\nwixident{Formula.dest{\_}iff}}{Formula.dest:uniff}}%
\nwixlogsorted{i}{{\nwixident{Formula.dest{\_}imp}}{Formula.dest:unimp}}%
\nwixlogsorted{i}{{\nwixident{Formula.dest{\_}not}}{Formula.dest:unnot}}%
\nwixlogsorted{i}{{\nwixident{Formula.dest{\_}or}}{Formula.dest:unor}}%
\nwixlogsorted{i}{{\nwixident{Formula.dest{\_}pred}}{Formula.dest:unpred}}%
\nwixlogsorted{i}{{\nwixident{Formula.fv}}{Formula.fv}}%
\nwixlogsorted{i}{{\nwixident{Formula.is{\_}and}}{Formula.is:unand}}%
\nwixlogsorted{i}{{\nwixident{Formula.is{\_}falsum}}{Formula.is:unfalsum}}%
\nwixlogsorted{i}{{\nwixident{Formula.is{\_}iff}}{Formula.is:uniff}}%
\nwixlogsorted{i}{{\nwixident{Formula.is{\_}imp}}{Formula.is:unimp}}%
\nwixlogsorted{i}{{\nwixident{Formula.is{\_}not}}{Formula.is:unnot}}%
\nwixlogsorted{i}{{\nwixident{Formula.is{\_}or}}{Formula.is:unor}}%
\nwixlogsorted{i}{{\nwixident{Formula.is{\_}verum}}{Formula.is:unverum}}%
\nwixlogsorted{i}{{\nwixident{Formula.mk{\_}and}}{Formula.mk:unand}}%
\nwixlogsorted{i}{{\nwixident{Formula.mk{\_}exists}}{Formula.mk:unexists}}%
\nwixlogsorted{i}{{\nwixident{Formula.mk{\_}falsum}}{Formula.mk:unfalsum}}%
\nwixlogsorted{i}{{\nwixident{Formula.mk{\_}forall}}{Formula.mk:unforall}}%
\nwixlogsorted{i}{{\nwixident{Formula.mk{\_}iff}}{Formula.mk:uniff}}%
\nwixlogsorted{i}{{\nwixident{Formula.mk{\_}imp}}{Formula.mk:unimp}}%
\nwixlogsorted{i}{{\nwixident{Formula.mk{\_}not}}{Formula.mk:unnot}}%
\nwixlogsorted{i}{{\nwixident{Formula.mk{\_}or}}{Formula.mk:unor}}%
\nwixlogsorted{i}{{\nwixident{Formula.mk{\_}pred}}{Formula.mk:unpred}}%
\nwixlogsorted{i}{{\nwixident{Formula.mk{\_}verum}}{Formula.mk:unverum}}%
\nwixlogsorted{i}{{\nwixident{Formula.subst}}{Formula.subst}}%
\nwixlogsorted{i}{{\nwixident{Formula.t}}{Formula.t}}%
\nwixlogsorted{i}{{\nwixident{Sort.IND}}{Sort.IND}}%
\nwixlogsorted{i}{{\nwixident{Sort.t}}{Sort.t}}%
\nwixlogsorted{i}{{\nwixident{Term.eq}}{Term.eq}}%
\nwixlogsorted{i}{{\nwixident{Term.sort}}{Term.sort}}%
\nwixlogsorted{i}{{\nwixident{Term.t}}{Term.t}}%
\nwixlogsorted{i}{{\nwixident{Unif.y{\_}complements}}{Unif.y:uncomplements}}%

